\chapter{Fundamental Theorem of Calculus}

\section{Exact Convexity}

\begin{definition}[Exact raw-real order]\label{def:raw-real-order}
\lean{ComputableAnalysis.RealRaw.Le}
\uses{def:raw-real}
\leanok
For raw reals \(X,Y\), define
\[
  X\le Y
\]
to mean that every rational lower approximation to \(X\) is at most every
rational upper approximation to \(Y\):
\[
  \forall n,m,\qquad X_n^- \le Y_m^+.
\]
This is not the executable one-stage comparison.  It is the exact order
relation used in mathematical statements.  For valid raw representatives,
two raw reals are equivalent exactly when each is at most the other in this
order.
\end{definition}

\begin{definition}[Endpoint difference real]\label{def:endpoint-diff}
\lean{ComputableAnalysis.endpointDifferenceInterval}
\uses{def:real-fun-raw,def:raw-real-arithmetic}
For a raw real-valued function \(F\), define
\[
  E(F;u,v)=F(v)-F(u).
\]
At stage \(p\), if \(F(u)_p=[A^-,A^+]\) and
\(F(v)_p=[B^-,B^+]\), then
\[
  E(F;u,v)_p=[B^- - A^+,B^+ - A^-].
\]
\end{definition}

\begin{definition}[Rational secant raw real]\label{def:rational-secant}
\lean{ComputableAnalysis.secantRaw}
\uses{def:endpoint-diff,def:raw-real-order}
\leanok
For rational \(u<v\), the secant of \(F\) on \([u,v]\) is the raw real
\[
  \Sec_F(u,v)=\frac{F(v)-F(u)}{v-u}.
\]
Its finite stages are the rational interval slopes
\[
  \Sec_F(u,v)_p=E(F;u,v)_p/(v-u).
\]
\end{definition}

\begin{definition}[Exact convexity]\label{def:exact-convexity}
\lean{ComputableAnalysis.ExactConvexOn}
\uses{def:rational-secant,def:raw-real-order}
\leanok
A raw real-valued function \(F\) is convex on the rational interval
\([a,b]\) when it is defined and valid at every rational point of
\([a,b]\), and its rational secants are monotone:
\[
  \Sec_F(w,x)\le \Sec_F(y,z)
\]
whenever
\[
  a\le w<x\le y<z\le b.
\]
Concavity is the same statement with the secant order reversed.
\end{definition}

\section{One-Sided Convex Derivatives}

\begin{definition}[Right derivative of a convex function]
\label{def:convex-right-derivative}
\lean{ComputableAnalysis.RightDerivativeAt}
\uses{def:exact-convexity}
Let \(F\) be convex on \([a,b]\), and let \(q<b\) be rational.  The right
secants
\[
  \Sec_F(q,q+h),\qquad h>0,\quad q+h\le b,
\]
are monotone as \(h\downarrow 0\) and bounded below by any left secant ending
at \(q\).  The right derivative \(D_+F(q)\) is the exact raw real represented
by the infimum of these right secants:
\[
  D_+F(q)=\inf_{0<h,\ q+h\le b}\Sec_F(q,q+h).
\]
Equivalently, it is the limit of any rational sequence \(h_n\downarrow0\)
with \(q+h_n\le b\).
\end{definition}

\begin{definition}[Left derivative of a convex function]
\label{def:convex-left-derivative}
\lean{ComputableAnalysis.LeftDerivativeAt}
\uses{def:exact-convexity}
For \(a<q\), the left derivative is
\[
  D_-F(q)=\sup_{0<h,\ a\le q-h}\Sec_F(q-h,q).
\]
For convex \(F\), \(D_-F(q)\le D_+F(q)\).  The ordinary two-sided derivative
exists exactly at the rational points where these two one-sided derivatives
are equivalent.
\end{definition}

\begin{theorem}[Right derivative is monotone]
\label{thm:convex-right-derivative-increasing}
\lean{ComputableAnalysis.rightDerivativeAt_mono}
\uses{def:convex-right-derivative}
\leanok
If \(F\) is convex on \([a,b]\), then
\[
  q_1<q_2
  \quad\Longrightarrow\quad
  D_+F(q_1)\le D_+F(q_2)
\]
for rational \(q_1,q_2<b\).
\end{theorem}

\begin{proof}
Choose rational \(h_1,h_2>0\) so small that
\[
  q_1+h_1\le q_2\le q_2+h_2.
\]
Convexity gives
\[
  \Sec_F(q_1,q_1+h_1)\le \Sec_F(q_2,q_2+h_2).
\]
Taking the infimum over right secants on both sides gives the claim.
\end{proof}

\begin{remark}[Corners]\label{rem:convex-corners}
\uses{def:convex-left-derivative,def:convex-right-derivative}
Convexity does not imply existence of the two-sided derivative.  The function
\(|x|\) is convex and has \(D_-F(0)<D_+F(0)\).  The universal FTC for convex
functions must therefore use a one-sided derivative, or else explicitly
restrict to the points where the left and right derivatives agree.
\end{remark}

\section{Integrating Monotone Derivatives}

\begin{definition}[Integral of a monotone raw function]
\label{def:monotone-integral}
\uses{def:raw-real-order}
Let \(g\) be a monotone rational-input raw real-valued function on
\([a,b]\).  Its integral is the common raw real obtained from lower and upper
Darboux sums over rational partitions:
\[
  L(P,g)\le \int_a^b g(x)\,dx \le U(P,g),
\]
with the mesh of \(P\) tending to zero.  Monotonicity supplies the cell
infimum and supremum from endpoint values, so the only analytic ingredient is
the completeness of raw reals for nested rational intervals.
\end{definition}

\begin{theorem}[Monotone functions are integrable]
\label{thm:monotone-integrable}
\uses{def:monotone-integral}
Every monotone rational-input raw real-valued function on a rational interval
has an integral in the sense of \cref{def:monotone-integral}.
\end{theorem}

\section{Convex FTC}

\begin{theorem}[Convex FTC]\label{thm:convex-ftc}
\uses{def:exact-convexity,def:convex-right-derivative,def:monotone-integral}
If \(F\) is convex on a rational interval \([a,b]\), then the right
derivative \(D_+F\) is integrable and its integral represents the same raw
real as \(F(b)-F(a)\).
Equivalently, the same conclusion holds with the left derivative \(D_-F\).
\end{theorem}

\begin{proof}
For a partition \(P:a=x_0<\cdots<x_N=b\), convexity gives
\[
  D_+F(x_i)
  \le
  \Sec_F(x_i,x_{i+1})
  \le
  D_-F(x_{i+1})
\]
on each cell.  Multiplying by \(x_{i+1}-x_i\) and summing yields Darboux
lower and upper sums enclosing
\[
  \sum_i \bigl(F(x_{i+1})-F(x_i)\bigr)=F(b)-F(a).
\]
The gap between the sums is controlled by monotonicity of the one-sided
derivatives, and tends to zero under rational refinement.  Hence the integral
and the endpoint difference represent the same raw real.
\end{proof}

\begin{remark}[No effective side conditions]\label{rem:no-effective-side-conditions}
\uses{thm:convex-ftc}
The theorem does not ask for derivative bounds, local controls, or overlap
estimates as hypotheses.  Those are proof artifacts from an earlier approach.
The mathematical input is exact convexity; the derivative, its monotonicity,
its integrability, and the endpoint equality are consequences.
\end{remark}

\begin{remark}[Piecewise convexity]\label{rem:no-piecewise-ftc}
\uses{thm:convex-ftc}
There is no separate named theorem for piecewise convex or concave functions.
In an example proof, split the interval at rational breakpoints where the
convexity proof changes, apply \cref{thm:convex-ftc} on each piece, and
combine the endpoint equalities by raw-real arithmetic and transitivity.
\end{remark}

\section{Formula Identification and Examples}

\begin{definition}[Formula identification]\label{def:formula-identification}
\uses{def:convex-right-derivative}
Let \(g\) be a proposed raw function formula.  On a rational interval,
\(g\) is identified with \(D_+F\) when it is equivalent to the right-secant
infimum defining \(D_+F\).  The formula is not a guess: it is proved from
finite secant inequalities or algebraic identities for the particular
function.
\end{definition}

\begin{theorem}[Formula identification closure]
\label{thm:formula-identification-ftc}
\uses{def:formula-identification,thm:convex-ftc}
If \(g\) is identified with \(D_+F\) on \([a,b]\), then
the integral of \(g\) represents the same raw real as \(F(b)-F(a)\).
\end{theorem}

\begin{example}[Polynomials]\label{ex:polynomial-ftc}
\uses{thm:formula-identification-ftc}
For rational polynomials, convexity or concavity on an interval is proved by
finite algebra, usually by certifying the sign of the second derivative on
that interval.  The derivative formula is the formal polynomial derivative.
\end{example}

\begin{example}[Rational functions]\label{ex:rational-function-ftc}
\uses{thm:formula-identification-ftc}
For rational functions, the domain must first be certified on the interval.
Denominator-apartness is a sufficient way to rule out hidden singularities
such as \(1/(x^2-2)\) on \([1,2]\).  After the domain is certified,
convexity and formula identification reduce to finite sign estimates.
\end{example}

\begin{example}[Power series]\label{ex:power-series-convex-ftc}
\uses{thm:formula-identification-ftc}
For a power series on a certified convergence interval, convexity can be
proved by enclosing the second differentiated series and its tail.  The
derivative formula is obtained by termwise differentiation, again with a tail
certificate.
\end{example}

\begin{example}[Geometric arctangent]\label{ex:arctan-ftc}
\uses{thm:formula-identification-ftc}
Let \(A(x)\) be the geometric arctangent, defined by sector-area intervals.
On \(0\le u<v\), the finite geometric inequality
\[
  \frac{1}{1+v^2}
  \le
  \Sec_A(u,v)
  \le
  \frac{1}{1+u^2}
\]
identifies the right derivative with
\[
  K(x)=\frac{1}{1+x^2}.
\]
The Taylor route then integrates the finite identity for \(K\) and its
explicit remainder.
\end{example}

\begin{example}[Logarithm]\label{ex:log-ftc}
\uses{thm:formula-identification-ftc}
For logarithm on \(0<a\le x\le b\), the derivative formula is \(1/x\).
The positivity certificate supplies the domain-apartness data needed for
interval division, and the secant inequalities identify \(1/x\) with the
one-sided convex derivative.
\end{example}

\begin{example}[Square root and arcsine]\label{ex:sqrt-arcsine-ftc}
\uses{thm:convex-ftc,thm:formula-identification-ftc}
The function \(\sqrt{x}\) is concave on positive intervals, so
\(-\sqrt{x}\) is convex.  The formula \(1/(2\sqrt{x})\) requires both a
certified square-root interval and apartness from \(0\).  For arcsine, the
kernel \((1-x^2)^{-1/2}\) similarly requires a certified interval inside
\((-1,1)\).
\end{example}
